\section{Text and image detector
catalogue}\label{record-DETECTORS-md-text-and-image-detector-catalogue}

This reference explains detector mechanisms, score interpretation, and
implementation requirements. For the research workflow, experiments, and
setup, see the \href{../README.md}{main README}.

\begin{itemize}
\tightlist
\item
  \hyperref[record-DETECTORS-md-how-the-text-detectors-work]{Text
  detectors}
\item
  \hyperref[record-DETECTORS-md-how-the-image-detectors-work]{Image
  detectors}
\end{itemize}

\subsection{How the text detectors
work}\label{record-DETECTORS-md-how-the-text-detectors-work}

\subsubsection{The example used
throughout}\label{record-DETECTORS-md-the-example-used-throughout}

The text examples use the following passage. All example scores, token
statistics, and model outputs in this catalogue are hypothetical
illustrations, not measurements from a recorded run. Model settings and
formulas describe the implementation.

Example passage:

\begin{quote}
Wild yeast and bacteria eat the sugars in the flour and give off carbon
dioxide, and the gluten network traps that gas. In the oven the starch
sets, so the loaf keeps the shape the gas gave it.
\end{quote}

Where a detector looks at individual tokens, the first four positions
are used: \texttt{\seqsplit{" Wild"}}, \texttt{\seqsplit{" yeast"}},
\texttt{\seqsplit{" and"}}, \texttt{\seqsplit{" bacteria"}}.

\subsubsection{Background: surprisal and
perplexity}\label{record-DETECTORS-md-background-surprisal-and-perplexity}

A language model predicts the next token from the preceding context. The
\textbf{surprise} (or surprisal) of an observed token \(x_t\) is
\(s_t=-\log p(x_t\mid x_{<t})\): a predictable token has low surprise.
For this passage the opening surprises are

{\scriptsize\setlength{\tabcolsep}{3pt}\def\LTcaptype{none} % do not increment counter
\begin{longtable}[]{@{}
  >{\raggedright\arraybackslash}p{(\linewidth - 2\tabcolsep) * \real{0.5000}}
  >{\raggedright\arraybackslash}p{(\linewidth - 2\tabcolsep) * \real{0.5000}}@{}}
\toprule\noalign{}
\begin{minipage}[b]{\linewidth}\raggedright
token
\end{minipage} & \begin{minipage}[b]{\linewidth}\raggedright
surprise
\end{minipage} \\
\midrule\noalign{}
\endhead
\bottomrule\noalign{}
\endlastfoot
\texttt{\seqsplit{" Wild"}} & 0.9 \\
\texttt{\seqsplit{" yeast"}} & 1.8 \\
\texttt{\seqsplit{" and"}} & 0.5 \\
\texttt{\seqsplit{" bacteria"}} & 1.6 \\
\end{longtable}
}

Perplexity is the exponential of the average surprise, so the average of
these four is about \texttt{\seqsplit{1.2}} nats, or perplexity
\texttt{\seqsplit{3.3}}. AI-generated text can have distinctive
predictability or variation patterns, but these are statistical
tendencies, not reliable signatures of authorship in every document.

\subsubsection{ModernBERT (supervised
baseline)}\label{record-DETECTORS-md-modernbert-supervised-baseline}

\texttt{\seqsplit{ModernBERTDetector}} learns which patterns distinguish
human and AI text from labeled examples. A QLoRA-fine-tuned ModernBERT
reads the passage and predicts its class; QLoRA trains compact adapters
on a mostly frozen, quantized model.

Suppose the classifier returns
\texttt{\seqsplit{details = \{"prob\_ai": 0.83, "prob\_human": 0.17\}}},
and the score is \texttt{\seqsplit{probs[0]}}, taken directly with no
transform.

\textbf{Interpretation:} the checkpoint's class order is
\texttt{\seqsplit{0 = machine}}, \texttt{\seqsplit{1 = human}}, so
\texttt{\seqsplit{0.83}} is the AI-oriented probability. A high score
means the passage resembles AI examples in the training distribution,
\textbf{not} that its authorship is known with that probability. Nothing
validates that class order, so a checkpoint trained the other way round
would be silently inverted. The input is truncated to the model's
\texttt{\seqsplit{max\_seq\_length}}, so later passages never reach the
classifier. Requires the \texttt{\seqsplit{bench\_aitextdetect}}
ModernBERT checkpoint.

\subsubsection{TF-IDF + logistic regression (supervised
baseline)}\label{record-DETECTORS-md-tf-idf-logistic-regression-supervised-baseline}

\texttt{\seqsplit{TFIDFDetector}} detects differences in word usage
rather than using a language model's token probabilities. TF-IDF
represents each document by words weighted more heavily when they are
frequent in that document but uncommon across the corpus. Logistic
regression learns which weighted words favor either class.

The weights that pushed this passage toward the AI side would be the
sort of tokens that separate the training classes, and the classifier
returns \texttt{\seqsplit{probs[0] = 0.62}}, again with no transform.

\textbf{Interpretation:} \texttt{\seqsplit{0.62}} sits in the
\texttt{\seqsplit{uncertain}} band. The signal is cheap and
interpretable, but it can change under paraphrasing or topic shift, and
classifier probabilities are not automatically well-calibrated. If
\texttt{\seqsplit{tfidf\_vectorizer.pkl}} is missing the detector
reports itself unavailable rather than fitting a fallback. Requires a
fitted vectorizer and classifier.

\subsubsection{Binoculars}\label{record-DETECTORS-md-binoculars}

\texttt{\seqsplit{BinocularsDetector}} detects AI-generated text by
comparing how predictable the observed tokens are to what a language
model would typically generate in the same context.

It uses two closely related language models:

\begin{itemize}
\tightlist
\item
  \textbf{Observer} (\(q_t\)): assigns probabilities to the actual
  tokens.
\item
  \textbf{Performer} (\(p_t\)): provides a reference distribution over
  possible tokens.
\end{itemize}

At each position, the numerator is the observer's surprise on the actual
token, and the denominator is the cross-entropy between the two
predictive distributions, i.e.~the observer's surprise averaged over the
performer's probabilities:

\[
B(x) =
\frac{\frac{1}{T}\sum_t -\log q_t(x_t)}
{\frac{1}{T}\sum_t \sum_v -p_t(v)\log q_t(v)}.
\]

{\scriptsize\setlength{\tabcolsep}{3pt}\def\LTcaptype{none} % do not increment counter
\begin{longtable}[]{@{}
  >{\raggedright\arraybackslash}p{(\linewidth - 4\tabcolsep) * \real{0.3333}}
  >{\raggedright\arraybackslash}p{(\linewidth - 4\tabcolsep) * \real{0.3333}}
  >{\raggedright\arraybackslash}p{(\linewidth - 4\tabcolsep) * \real{0.3333}}@{}}
\toprule\noalign{}
\begin{minipage}[b]{\linewidth}\raggedright
position
\end{minipage} & \begin{minipage}[b]{\linewidth}\raggedright
actual surprise \(-\log q_t(x_t)\)
\end{minipage} & \begin{minipage}[b]{\linewidth}\raggedright
expected \(\sum_v -p_t(v)\log q_t(v)\)
\end{minipage} \\
\midrule\noalign{}
\endhead
\bottomrule\noalign{}
\endlastfoot
\texttt{\seqsplit{" Wild"}} & 0.9 & 1.4 \\
\texttt{\seqsplit{" yeast"}} & 1.8 & 2.2 \\
\texttt{\seqsplit{" and"}} & 0.5 & 1.1 \\
\texttt{\seqsplit{" bacteria"}} & 1.6 & 2.0 \\
mean & 1.20 & 1.68 \\
\end{longtable}
}

So \(B(x) = 1.20 / 1.68 \approx 0.72\).

\textbf{Interpretation:} \textbf{lower is more AI.} The observed tokens
were more predictable than the performer's reference distribution
expected, which is the hypothesis the method encodes. The default
threshold is \texttt{\seqsplit{0.9015}}: \(B(x) < 0.9015\) is labeled
\textbf{AI-generated}, \(B(x) \geq 0.9015\) is labeled
\textbf{human-written}, and the label is taken from the raw ratio, so it
is always binary. This is a low-false-positive-rate operating point, not
a guarantee. The display score is
\texttt{\seqsplit{sigmoid((0.9015 - raw) / 0.05)}}, which maps
\texttt{\seqsplit{0.72}} to about \texttt{\seqsplit{0.97}} and the
threshold itself to \texttt{\seqsplit{0.5}}. A calibrator can replace
that map; the raw ratio and the threshold used are kept in
\texttt{\seqsplit{details}}.

Both models are truncated to \texttt{\seqsplit{512}} tokens, and the
observer and performer are not interchangeable: swapping them changes
the statistic. The defaults are \texttt{\seqsplit{tiiuae/falcon-7b}}
(observer) and \texttt{\seqsplit{tiiuae/falcon-7b-instruct}}
(performer), both requiring CUDA.

\subsubsection{Fast-DetectGPT}\label{record-DETECTORS-md-fast-detectgpt}

\texttt{\seqsplit{FastDetectGPTDetector}} asks whether the observed
tokens are unusually likely under a scoring language model, relative to
plausible alternative tokens. It compares each actual token's
log-probability with the expected log-probability of tokens drawn from a
reference model, then normalizes the difference by its expected
variation. Unlike perturbation-based methods, it does \textbf{not} need
to generate rewritten versions of the text.

Let \(\ell_t(v)\) be the scoring model's log-probability and \(r_t(v)\)
the reference distribution at position \(t\):

\[
C(x)=\frac{\sum_t[\ell_t(x_t)-\mathbb E_{v\sim r_t}\ell_t(v)]}
{\sqrt{\sum_t\operatorname{Var}_{v\sim r_t}[\ell_t(v)]}}.
\]

{\scriptsize\setlength{\tabcolsep}{3pt}\def\LTcaptype{none} % do not increment counter
\begin{longtable}[]{@{}
  >{\raggedright\arraybackslash}p{(\linewidth - 8\tabcolsep) * \real{0.2000}}
  >{\raggedright\arraybackslash}p{(\linewidth - 8\tabcolsep) * \real{0.2000}}
  >{\raggedright\arraybackslash}p{(\linewidth - 8\tabcolsep) * \real{0.2000}}
  >{\raggedright\arraybackslash}p{(\linewidth - 8\tabcolsep) * \real{0.2000}}
  >{\raggedright\arraybackslash}p{(\linewidth - 8\tabcolsep) * \real{0.2000}}@{}}
\toprule\noalign{}
\begin{minipage}[b]{\linewidth}\raggedright
position
\end{minipage} & \begin{minipage}[b]{\linewidth}\raggedright
\(\ell_t(x_t)\)
\end{minipage} & \begin{minipage}[b]{\linewidth}\raggedright
\(\mathbb E_{v\sim r_t}\ell_t(v)\)
\end{minipage} & \begin{minipage}[b]{\linewidth}\raggedright
\(\operatorname{Var}_{v\sim r_t}\ell_t\)
\end{minipage} & \begin{minipage}[b]{\linewidth}\raggedright
difference
\end{minipage} \\
\midrule\noalign{}
\endhead
\bottomrule\noalign{}
\endlastfoot
\texttt{\seqsplit{" Wild"}} & -0.9 & -1.4 & 0.09 & +0.5 \\
\texttt{\seqsplit{" yeast"}} & -1.8 & -2.2 & 0.09 & +0.4 \\
\texttt{\seqsplit{" and"}} & -0.5 & -1.1 & 0.16 & +0.6 \\
\texttt{\seqsplit{" bacteria"}} & -1.6 & -2.0 & 0.16 & +0.4 \\
\end{longtable}
}

So \(C(x) = 1.9 / \sqrt{0.50} \approx 2.69\).

\textbf{Interpretation:} a larger \(C(x)\) means the actual tokens are
more favored than typical reference alternatives, relative to the
variance; this is treated as \textbf{more AI-like}. Classification
compares the raw statistic to its configured threshold, default
\texttt{\seqsplit{0.0}}, so the label is binary and both values are kept
in \texttt{\seqsplit{details}}. The \texttt{\seqsplit{[0, 1]}} display
score is a plain \texttt{\seqsplit{sigmoid(C(x))}}, here about
\texttt{\seqsplit{0.94}}, and \texttt{\seqsplit{score\_semantics}} is
\texttt{\seqsplit{uncalibrated\_logistic\_transform}}.

Two implementation details matter when reading raw values. The numerator
and the variance are both \textbf{sums}, not means, so the statistic
grows roughly like \(\sqrt{N}\) with length and is not comparable across
passage sizes. And by default the reference model \textbf{is} the
scoring model, which reduces the expectation to the negated entropy of
that model; pass an explicit second model to get the paper's
construction. When you do, the two tokenizers must have identical
vocabularies. Inputs are truncated to \texttt{\seqsplit{512}} tokens,
and passages shorter than ten tokens return a score of
\texttt{\seqsplit{0.0}}.

\subsubsection{DivEye}\label{record-DETECTORS-md-diveye}

\texttt{\seqsplit{DivEyeDetector}} now uses the authors' released
\textbf{combined DivEye/BiScope} classifier, rather than the local
standalone nine-feature approximation. It concatenates eleven Falcon-7B
surprisal/compression features with 72 Gemma-1.1- 2B-Instruct BiScope
features. XGBoost maps that ordered 83-feature vector to an AI
probability, using probability column 1 as in the authors' application.

Both language models and the already fitted classifier remain frozen
during inference and attacker training. No local classifier fitting is
required. The implementation pins the classifier revision and checksum
and both model revisions, rejects incompatible classifier schemas, and
supports one GPU. Inputs require at least 15 words; Falcon truncates at
1,024 tokens and Gemma at 2,000 text tokens plus the completion prompt.
Gemma requires authorized Hugging Face access. Detection details
identify the combined variant, classifier hash, model identities and
score orientation. This release is not the paper's standalone
nine-feature experiment. See
\href{records/docs/diveye_classifier.tex}{integration notes}.

\subsubsection{Disrupt-and-Recover}\label{record-DETECTORS-md-disrupt-and-recover}

\texttt{\seqsplit{DisruptRecoverDetector}} tests how easily a model can
reconstruct the original word order after scrambling it. The hypothesis
is that machine-written prose may have more regular structure, making
its original arrangement easier to recover.

It splits the text at punctuation boundaries and shuffles the words
within each chunk of at least three words, using a fixed seed so the
corruption is reproducible. For this passage the second chunk becomes:

\begin{quote}
that gas the gluten network traps.
\end{quote}

The recovery model is asked to restore the order without adding or
removing words, giving back
\texttt{\seqsplit{the gluten network traps that gas.}} The result is
compared with the original in two ways: \textbf{BERTScore} F1
(\texttt{\seqsplit{0.94}}) for semantic similarity, and
\texttt{\seqsplit{coverage x mean(Kendall, Spearman)}} remapped to
\texttt{\seqsplit{[0, 1]}} (\texttt{\seqsplit{0.83}}) for word order
agreement. The two are averaged to give a
\texttt{\seqsplit{recovery\_similarity}} of \texttt{\seqsplit{0.885}}.

\textbf{Interpretation:} greater recovery similarity maps to a more
AI-oriented score, here
\texttt{\seqsplit{sigmoid((0.885 - 0.72) / 0.08) = 0.89}}. The
\texttt{\seqsplit{0.72}} is the \textbf{center of that logistic}, not a
decision boundary: the \texttt{\seqsplit{0.65}} label band is only
reached at a similarity of about \texttt{\seqsplit{0.77}}. This is a
heuristic hypothesis, not a guaranteed property of AI text, and no
probability interpretation follows from the score. Passages under 35
words return \texttt{\seqsplit{uncertain}} with a reason. Requires an
injected \texttt{\seqsplit{recover\_fn}} or an OpenAI-compatible model
configured with \texttt{\seqsplit{DR\_RECOVERY\_MODEL}} and an API key,
plus BERTScore weights. The corrupted and recovered texts are stored in
\texttt{\seqsplit{details}}, which makes results large.

\subsubsection{Window
smoothing}\label{record-DETECTORS-md-window-smoothing}

A \textbf{word window} is a consecutive chunk of text. With
\texttt{\seqsplit{window\_size = 64}} and
\texttt{\seqsplit{stride = 32}}, the detector reads words 1-64, then
33-96, then 65-128, and so on. The starting position moves forward by 32
words each time, so adjacent windows share 32 words. Here words are
separated by whitespace, not tokenizer tokens.

\texttt{\seqsplit{WindowSmoothedTextDetector}} runs an existing detector
on each chunk, adjusts its score using neighboring chunks' scores, then
combines them into one document score. The purpose is to examine
passages individually and reduce isolated score fluctuations: if one
chunk scores strongly AI-like while its overlapping neighbors score
human-like, their scores influence its adjustment. Texts of 64 words or
fewer are scored as a single window.

This is a heuristic wrapper around a detector, not new evidence of
authorship. It can also blur a real change between human and AI
passages, and improved accuracy must be validated. The first two windows
receive reduced weights inside the adjustment and final average; an
optional complexity penalty reduces the score for very simple text.
Overlapping words contribute more than once, so this is not a per-word
average. Raw and adjusted scores are returned in
\texttt{\seqsplit{details}}. This implementation is not the cited
paper's learned token-level Markov random field;
\texttt{\seqsplit{MarkovCalibratedTextDetector}} is a compatibility
alias.

\subsubsection{IPAD}\label{record-DETECTORS-md-ipad}

\texttt{\seqsplit{IPADDetector}} asks whether a passage looks like the
answer to a plausible AI prompt. Its key idea is to \textbf{infer a
prompt from the text}, then check how well the prompt and the text fit
together, including via an independently regenerated answer.

Suppose the inverter adapter returns the prompt
\texttt{\seqsplit{Explain why sourdough bread rises.}}, and the other
two adapters judge that prompt against the text:

{\scriptsize\setlength{\tabcolsep}{3pt}\def\LTcaptype{none} % do not increment counter
\begin{longtable}[]{@{}
  >{\raggedright\arraybackslash}p{(\linewidth - 4\tabcolsep) * \real{0.3333}}
  >{\raggedright\arraybackslash}p{(\linewidth - 4\tabcolsep) * \real{0.3333}}
  >{\raggedright\arraybackslash}p{(\linewidth - 4\tabcolsep) * \real{0.3333}}@{}}
\toprule\noalign{}
\begin{minipage}[b]{\linewidth}\raggedright
adapter
\end{minipage} & \begin{minipage}[b]{\linewidth}\raggedright
question
\end{minipage} & \begin{minipage}[b]{\linewidth}\raggedright
answer probability
\end{minipage} \\
\midrule\noalign{}
\endhead
\bottomrule\noalign{}
\endlastfoot
PTCV & could an AI have generated this text from that prompt? &
\texttt{\seqsplit{p\_yes = 0.80}} \\
RC & is this text generated similarly to a fresh answer to that prompt?
& \texttt{\seqsplit{p\_yes = 0.60}} \\
\end{longtable}
}

Both are two-way renormalized sequence likelihoods over the tokens
\texttt{\seqsplit{" yes"}} and \texttt{\seqsplit{" no"}}, so they sum to
\texttt{\seqsplit{1}} by construction. The fused score is
\texttt{\seqsplit{0.45 * 0.80 + 0.55 * 0.60 = 0.69}}.

\textbf{Interpretation:} the paper's decision threshold is
\texttt{\seqsplit{0.54}} and the comparison is strict, so
\texttt{\seqsplit{0.69}} is labeled \textbf{AI}. A higher fused score
indicates greater consistency with the inferred-prompt generation
process. This is a method-specific threshold, not universal proof or a
guaranteed calibrated probability. The fusion gives PTCV weight
\texttt{\seqsplit{0.45}} and RC weight \texttt{\seqsplit{0.55}}.
Requires the approximately 14B-parameter base model, all three adapters,
and a sufficiently large GPU; CPU loading is refused. Regeneration is a
separate operation, not a fourth adapter, and
\texttt{\seqsplit{details["regeneration\_backend"]}} records which model
produced it.

\subsubsection{Paraphrase round
trip}\label{record-DETECTORS-md-paraphrase-round-trip}

\texttt{\seqsplit{ParaphraseRoundTripDetector}} checks \textbf{how far
the passage is from plain, predictable prose}. It asks a generator to
rewrite the text ``into clear, standard, and predictable prose'', then
measures the \textbf{word-level} edit distance between the original and
the rewrite, normalized by the longer of the two.

If the rewrite changed 12 of the passage's 34 words, the normalized
distance is \texttt{\seqsplit{0.35}}, and the score is
\texttt{\seqsplit{min(0.35 * 1.5, 1.0) = 0.53}}.

\textbf{Interpretation:} the mapping is a linear ramp with gain
\texttt{\seqsplit{1.5}} that saturates at a normalized distance of
\texttt{\seqsplit{0.67}}, so it is a heuristic normalized distance, not
a semantic similarity and \textbf{not} a calibrated probability; an
optional complexity penalty multiplies the score by
\texttt{\seqsplit{0.8}} for very simple text. The direction is the
stated one: a large distance means the original was far from plain
prose, which is scored as AI. This is \textbf{not} a validated universal
authorship test, and the 20-character minimum is a character count
rather than the 35-word minimum used by Disrupt-and-Recover. The full
rewrite is stored in \texttt{\seqsplit{details["rewritten"]}}.
\texttt{\seqsplit{InversionDetector}} is a compatibility alias.

\subsubsection{LUAR style
embedding}\label{record-DETECTORS-md-luar-style-embedding}

\texttt{\seqsplit{StyleEmbeddingDetector}} asks whether a document's
writing \textbf{style} more closely resembles known human or AI
examples. LUAR maps texts into an embedding space; the method computes
human and AI support-set centroids and compares the query embedding with
each using cosine similarity.

For this passage, with a support set of 200 AI and 200 human texts:

{\scriptsize\setlength{\tabcolsep}{3pt}\def\LTcaptype{none} % do not increment counter
\begin{longtable}[]{@{}
  >{\raggedright\arraybackslash}p{(\linewidth - 2\tabcolsep) * \real{0.5000}}
  >{\raggedright\arraybackslash}p{(\linewidth - 2\tabcolsep) * \real{0.5000}}@{}}
\toprule\noalign{}
\begin{minipage}[b]{\linewidth}\raggedright
quantity
\end{minipage} & \begin{minipage}[b]{\linewidth}\raggedright
value
\end{minipage} \\
\midrule\noalign{}
\endhead
\bottomrule\noalign{}
\endlastfoot
cosine to AI centroid & 0.71 \\
cosine to human centroid & 0.33 \\
\end{longtable}
}

The score is a softmax over those two cosines with temperature
\texttt{\seqsplit{0.5}}, giving
\texttt{\seqsplit{exp(0.71/0.5) / (exp(0.71/0.5) + exp(0.33/0.5)) = 0.68}},
not the raw cosine.

\textbf{Interpretation:} a higher AI-oriented score indicates greater
stylistic similarity to the supplied \textbf{AI support examples},
relative to the human ones. It depends critically on the support set,
not on a universal pretrained ``AI style'' prototype, and the
temperature decides how sharp the boundary is. Both support classes must
be provided with \texttt{\seqsplit{setup\_support\_set()}} before
detection; the detector will not fabricate a default.

\subsection{How the image detectors
work}\label{record-DETECTORS-md-how-the-image-detectors-work}

\subsubsection{DDA and Community
Forensics}\label{record-DETECTORS-md-dda-and-community-forensics}

\texttt{\seqsplit{DDADetector}} loads the released NeurIPS 2025 DDA
classifier: DINOv2 ViT-L/14 with rank-8 LoRA and a binary head. The
official checkpoint and DINOv2/DDA source revisions are pinned in the
image experiment resource manifest. Native inference uses a 336-pixel
center crop with zero padding for smaller images and CLIP normalization.
Its sigmoid output is oriented \texttt{\seqsplit{1 = fake}}. The SD1.5
race explicitly resizes both classes to 200 pixels before that crop.

\texttt{\seqsplit{CommunityForensicsDetector}} loads the authors'
released CVPR 2025 ViT-S checkpoints through \texttt{\seqsplit{timm}},
with strict state loading. The 224 model resizes the short side to 256;
the 384 model resizes it to 440, then center crops to the model's input
size and uses ImageNet normalization. The sigmoid output is
\texttt{\seqsplit{1 = fake}}. Revisions are pinned, availability checks
are local-only, and model weights load lazily. This detector is not used
in the current experiment.

Both return classifier probabilities, without claiming calibration or
resistance to adaptive attacks. See the
\href{records/docs/image_defenders.tex}{paper review},
\href{https://github.com/roy-ch/Dual-Data-Alignment}{DDA code}, and
\href{https://github.com/JeongsooP/Community-Forensics}{Community
Forensics code}.

\subsubsection{EfficientNet (supervised
baseline)}\label{record-DETECTORS-md-efficientnet-supervised-baseline}

\texttt{\seqsplit{EfficientNetDetector}} uses convolutional features to
classify real and generated images. The classifier is trained on labeled
examples; its spatial filters can capture textures and artifacts
associated with generators in the training data. This wrapper loads the
trained checkpoint rather than training it during scoring.

PIL inputs are resized to \texttt{\seqsplit{200 x 200}} and
ImageNet-normalized. The score is the classifier's probability for class
\texttt{\seqsplit{0 = fake}}: higher means more resemblance to generated
examples learned during training, not proof of origin. The wrapper
assumes that class order without validating it. Requires the internal
\texttt{\seqsplit{bench\_imai\_artifact}} EfficientNet checkpoint.

\subsubsection{CLIP image classifier (supervised
baseline)}\label{record-DETECTORS-md-clip-image-classifier-supervised-baseline}

\texttt{\seqsplit{CLIPImageDetector}} uses pretrained CLIP image
representations with a binary classifier trained on labeled real and
generated images. The fitted classifier learns which features
distinguish the classes; CLIP alone does not supply an AI-authorship
decision. This wrapper loads the trained checkpoint, rather than
comparing the image with text prompts as in zero-shot CLIP.

It uses the same preprocessing and scoring convention as EfficientNet.
The score is the classifier's probability for class
\texttt{\seqsplit{0 = fake}}, so higher means more AI-like according to
the fitted classifier. The assumed class order must match the
checkpoint. Requires the internal
\texttt{\seqsplit{bench\_imai\_artifact}} CLIP checkpoint.

\subsubsection{DINOv2 image classifier (supervised
baseline)}\label{record-DETECTORS-md-dinov2-image-classifier-supervised-baseline}

\texttt{\seqsplit{DINOv2Detector}} uses pretrained, self-supervised
DINOv2 image features with a binary classifier trained on labeled real
and generated images. DINOv2 provides visual representations; the fitted
classifier learns which patterns favor the AI class. This wrapper loads
the trained checkpoint.

It uses the same preprocessing and class convention as EfficientNet: the
score is the probability for class \texttt{\seqsplit{0 = fake}}, with
higher values meaning more AI-like according to the fitted classifier.
Without that classifier, DINOv2 alone does not assign authorship labels.
This differs from \texttt{\seqsplit{WaRPADDetector}}, which compares
DINOv2 representations without fitting an authorship classifier.
Requires the internal \texttt{\seqsplit{bench\_imai\_artifact}} DINOv2
checkpoint.

\subsubsection{Patch-based
classifier}\label{record-DETECTORS-md-patch-based-classifier}

\texttt{\seqsplit{PatchBasedClassifier}} scores local image regions with
an existing detector. It slides \texttt{\seqsplit{128 x 128}} crops with
stride \texttt{\seqsplit{64}} over the image, then returns the
\textbf{85th percentile} of their scores. This emphasizes high-scoring
regions rather than averaging their evidence away across the whole
image.

No additional model is trained. The result depends on the underlying
detector, crop size, and number of patches; improved detection must be
validated. If either image dimension is at or below the patch size, the
wrapper scores the whole image instead. Requires a functioning
underlying detector.

\subsubsection{Frequency information and Haar
wavelets}\label{record-DETECTORS-md-frequency-information-and-haar-wavelets}

\texttt{\seqsplit{FrequencyDetector}} asks whether an image contains an
unusual balance of fine detail and coarse structure. The image is
converted to grayscale at native resolution and a one-level Haar
transform separates low-frequency information (broad color and shape,
\texttt{\seqsplit{LL}}) from high-frequency information (edges and
texture, \texttt{\seqsplit{LH + HL + HH}}). The detector measures the
\textbf{high-to-low energy ratio} and maps it with a logistic of
steepness \texttt{\seqsplit{20}} centered on \texttt{\seqsplit{0.15}}:

\[
S = \frac{1}{1 + \exp\left(-20\left(\frac{E_{LH} + E_{HL} + E_{HH}}{E_{LL}} - 0.15\right)\right)}.
\]

For an illustrative ratio of \texttt{\seqsplit{0.19}}, \(S = 0.69\).

\textbf{Interpretation:} the sigmoid is extremely steep, so
\texttt{\seqsplit{0.05}} maps to \texttt{\seqsplit{0.12}} and
\texttt{\seqsplit{0.35}} maps to \texttt{\seqsplit{0.98}}; the score
saturates well before the ratio does. The center
\texttt{\seqsplit{0.15}} is a hand-chosen midpoint, not a fitted
threshold and not the \texttt{\seqsplit{0.65}} label band. The statistic
is \textbf{relative} high-frequency energy and is computed at native
resolution, so natural texture, resizing, sharpening, and compression
all move it; grayscale conversion discards chromatic artifacts entirely.
This cheap heuristic is \textbf{not} WaRPAD and its score is not a
calibrated probability.

\subsubsection{MLEP}\label{record-DETECTORS-md-mlep}

\texttt{\seqsplit{MLEPDetector}} searches for \textbf{local texture
irregularities across image scales}. It independently shuffles
non-overlapping \texttt{\seqsplit{2 x 2}} patches in each color channel,
downscales and restores each of the scales \texttt{\seqsplit{1}},
\texttt{\seqsplit{0.5}}, \texttt{\seqsplit{0.25}}, and builds exact
stride-one \texttt{\seqsplit{2 x 2}} local entropy maps, which are
stacked into a \texttt{\seqsplit{(255, 255, 9)}} map and scaled by
\texttt{\seqsplit{2}} into \texttt{\seqsplit{[0, 1]}} before a ResNet-50
classifier.

For an illustrative classifier logit of \texttt{\seqsplit{1.10}}, the
score is \texttt{\seqsplit{sigmoid(1.10) = 0.75}}.

\textbf{Interpretation:} a higher output means the multiscale entropy
patterns of a \textbf{deliberately shuffled} image resemble generated
images in the classifier's training data. The shuffling happens before
the resampling, so the statistic describes a destroyed local structure,
and reproducing a number requires the same shuffle seed. The permutation
and entropy computation do \textbf{not} themselves establish a fake
label: the fitted CNN checkpoint is mandatory, its
\texttt{\seqsplit{conv1}} must accept
\texttt{\seqsplit{3 * len(scales)}} channels, and the load is strict.
Changing \texttt{\seqsplit{scales}} invalidates a previously trained
checkpoint.

\subsubsection{WaRPAD}\label{record-DETECTORS-md-warpad}

\texttt{\seqsplit{WaRPADDetector}} tests whether removing fine image
detail changes the image's \textbf{high-level representation unusually
strongly}. It resizes to \texttt{\seqsplit{896 x 896}}, removes
high-frequency content with a two-level Haar transform
(\texttt{\seqsplit{x - 0.1 * HF}}), tiles the result into
\texttt{\seqsplit{4 x 4 = 16}} patches of \texttt{\seqsplit{224 x 224}},
and compares DINOv2 ViT-L/14 CLS embeddings of the original and
perturbed patches.

The score is one minus the mean cosine similarity, clipped to
\texttt{\seqsplit{[0, 1]}}. For an illustrative mean cosine of
\texttt{\seqsplit{0.82}}, this gives
\texttt{\seqsplit{1 - 0.82 = 0.18}}.

\textbf{Interpretation:} a high cosine indicates representations that
are robust to detail removal, which the method treats as more real-like.
There is no sigmoid and no fitted threshold; with neither a
\texttt{\seqsplit{threshold}} nor a \texttt{\seqsplit{calibrator}}
supplied, the default label is \texttt{\seqsplit{"uncertain"}} for
\textbf{every} image, so the \texttt{\seqsplit{0.18}} above is a usable
number with no decision attached. Supply a threshold and the label is
taken from the raw score; supply a calibrator and
\texttt{\seqsplit{score\_semantics}} becomes
\texttt{\seqsplit{calibrated\_probability}}. Requires DINOv2 weights.
One reproduction detail: the Hugging Face processor resizes each
already-\texttt{\seqsplit{224}} patch to \texttt{\seqsplit{256}} and
center-crops back, so the outer margin of every patch never reaches the
encoder.

\subsubsection{Denoising trajectory
detection}\label{record-DETECTORS-md-denoising-trajectory-detection}

\texttt{\seqsplit{DenoisingTrajectoryDetector}} asks whether real and
generated images behave differently when moved backward and forward
along a diffusion process. It encodes the image into a Stable Diffusion
latent using the distribution \textbf{mean} (deterministic), applies
DDIM inversion in ascending \(t\) over \texttt{\seqsplit{50}} steps,
decodes each predicted-clean latent, and compares CLIP ViT-L/14
layer-\texttt{\seqsplit{15}} CLS features against the original.

For an illustrative mean cosine similarity of \texttt{\seqsplit{0.41}}
over the \texttt{\seqsplit{50}} steps, the score is

\[
S = \mathrm{clip}\left(\frac{0.41 + 1}{2}, 0, 1\right) = 0.705,
\]

which crosses the \texttt{\seqsplit{0.65}} band in cosine space at a
similarity of \texttt{\seqsplit{0.30}}.

\textbf{Interpretation:} trajectory-level feature consistency is the
detection signal, but read the direction carefully: the implemented
score \textbf{increases} with reconstruction fidelity, so a higher value
means the image reconstructs well through inversion. The returned value
is an \textbf{uncalibrated similarity transform} with
\texttt{\seqsplit{score\_semantics}} set accordingly, and the threshold
is a fixed affine map that was never fitted, so many images land in
\texttt{\seqsplit{uncertain}}. Exact numerical parity with the authors'
code and weights has \textbf{not} been established. Cost is 50 UNet
forwards plus 50 sequential VAE decodes per image, and
\texttt{\seqsplit{summed\_clip\_similarity}} in
\texttt{\seqsplit{details}} scales with the number of selected steps, so
it is not comparable across \texttt{\seqsplit{trajectory\_indices}}
settings. Requires the diffusion and CLIP models.
